\section{Electromagnetic closure of the optical interface}
\label{sec:r13-em-laws}

\begin{panel}
R12 defines $U_{\mathrm{det},\lambda}=H_{z,\lambda}[A\,t_\lambda
e^{i\phi_{\mathrm{lens},\lambda}}U_{\mathrm{in},\lambda}]$.
This chapter supplies a physical route to its transfer, loss and material-state
arguments. Maxwell's equations are the adopted classical field laws; a material
law, geometry and boundary data select a particular device. The added
templates fit the existing exact expression and finite-state interfaces.
\end{panel}

\subsection{Fields, charge and the data that close their evolution}
\label{sec:r13-em-maxwell}

Use SI quantities in one declared spatial frame and time coordinate. For a
stationary macroscopic medium, with free charge $\rho_f$ and free current
$\mathbf J_f$, adopt
\begin{align}
 \nabla\cdot\mathbf D&=\rho_f,& \nabla\cdot\mathbf B&=0,\nonumber\\
 \nabla\times\mathbf E&=-\partial_t\mathbf B,&
 \nabla\times\mathbf H&=\mathbf J_f+\partial_t\mathbf D.
 \label{eq:r13-em-maxwell}
\end{align}
The dimensions are $[\mathbf E]=\mathrm{V/m}$,
$[\mathbf H]=\mathrm{A/m}$, $[\mathbf D]=\mathrm{C/m^2}$,
$[\mathbf B]=\mathrm{T}$, $[\rho_f]=\mathrm{C/m^3}$ and
$[\mathbf J_f]=\mathrm{A/m^2}$. The divergence of the last equation gives
\begin{equation}
 \partial_t\rho_f+\nabla\cdot\mathbf J_f=0.
 \label{eq:r13-em-charge}
\end{equation}
Thus supplied charge and current histories cannot be chosen independently.
The electric divergence constraint propagates when it holds initially and
the supplied charge obeys continuity; the magnetic constraint propagates
because the divergence of a curl vanishes. These statements require the
corresponding classical derivatives or an explicitly declared weak solution.
The laws and charge-conservation argument are documented in the
\href{https://www.feynmanlectures.caltech.edu/II_18.html}{Caltech Feynman Lectures, II.18 [EM1]}.

For a stationary interface, let $\mathbf n$ point from medium 1 to medium 2.
With surface free charge $\rho_s$, surface free current $\mathbf K_s$, and
no magnetic surface sources, the interface conditions are
\begin{align}
 \mathbf n\cdot(\mathbf D_2-\mathbf D_1)&=\rho_s,&
 \mathbf n\cdot(\mathbf B_2-\mathbf B_1)&=0,\nonumber\\
 \mathbf n\times(\mathbf E_2-\mathbf E_1)&=0,&
 \mathbf n\times(\mathbf H_2-\mathbf H_1)&=\mathbf K_s.
 \label{eq:r13-em-boundary}
\end{align}
They follow by shrinking Maxwell's integral pillboxes and loops across the
interface. A moving boundary needs the corresponding moving-medium model;
these stationary formulas cannot silently serve as that model.
\href{https://live.ocw.mit.edu/courses/6-013-electromagnetics-and-applications-spring-2009/d3be4ea78b036a6362230fb41780cf54_MIT6_013S09_notes.pdf}{MIT Electromagnetics and Applications, Section 2.6 [EM2]}.

A device profile binds its material regions, interface geometry, constitutive
relations, sources, compatible initial fields and initial material memory.
It selects the needed independent boundary or incoming-wave data, including
an outgoing-wave condition for an open problem. It does not prescribe every
field component independently on every boundary. Existence and uniqueness
belong to this completed problem; the four field equations alone do not
resolve a resonant, contradictory or underspecified setup.

\subsection{Power balance before optical loss is called heat}
\label{sec:r13-em-energy}

The vector identity for $\nabla\cdot(\mathbf E\times\mathbf H)$ and
equation~\eqref{eq:r13-em-maxwell} give the constitutive-independent identity
\begin{equation}
 \nabla\cdot\mathbf S+\mathbf E\cdot\partial_t\mathbf D
 +\mathbf H\cdot\partial_t\mathbf B=-\mathbf E\cdot\mathbf J_f,
 \qquad \mathbf S=\mathbf E\times\mathbf H.
 \label{eq:r13-em-poynting-general}
\end{equation}
For instantaneous, time-independent real symmetric positive-definite tensors
$\epsilon,\mu$, with $\mathbf D=\epsilon\mathbf E$ and
$\mathbf B=\mu\mathbf H$, set
\begin{align}
 u_{\rm em}&=\tfrac12\mathbf E\cdot\epsilon\mathbf E+
             \tfrac12\mathbf H\cdot\mu\mathbf H,\nonumber\\
 \partial_tu_{\rm em}+\nabla\cdot\mathbf S
             &=-\mathbf E\cdot\mathbf J_f.
 \label{eq:r13-em-poynting-simple}
\end{align}
The quadratic energy follows by differentiating each form; time independence
and symmetry are used at that exact step. This is Poynting's energy balance,
with fields doing work on matter through $\mathbf E\cdot\mathbf J_f$.
\href{https://www.feynmanlectures.caltech.edu/II_27.html}{Caltech Feynman Lectures, II.27 [EM3]}.

Writing $\mathbf J_f=\sigma\mathbf E+\mathbf J_{\rm src}$, where the symmetric
part of $\sigma$ is positive semidefinite, identifies a nonnegative Ohmic sink
$q_{\rm Ohm}=\mathbf E\cdot\sigma\mathbf E$. The impressed source term
$-\mathbf E\cdot\mathbf J_{\rm src}$ may supply field energy. In contrast,
arbitrary $\mathbf E\cdot\mathbf J_f$ need not be positive or entirely heat:
work can enter mechanical or other material stores. A dispersive dielectric
also cannot use $\tfrac12\mathbf E\cdot\mathbf D$ as its complete instantaneous
energy without a justified material-state model.

\subsection{A finite dispersive material law with an energy proof}
\label{sec:r13-em-lorentz}

An optional, explicit local linear closure uses finitely many polarization
oscillators. In a fixed material region let $\epsilon_b,\mu>0$ be scalar,
time-independent background coefficients and set
\begin{align}
 \mathbf D&=\epsilon_b\mathbf E+\sum_{j=1}^{m}\mathbf P_j,
 &\mathbf B&=\mu\mathbf H,\nonumber\\
 \dot{\mathbf P}_j&=\mathbf V_j,
 &\dot{\mathbf V}_j&=f_j\mathbf E-\gamma_j\mathbf V_j-
                              \omega_j^2\mathbf P_j,
 \label{eq:r13-em-lorentz}
\end{align}
where $m$ is finite, $f_j>0$, $\gamma_j\ge0$ and $\omega_j>0$.
Here $[f_j]=[\epsilon_b]/\mathrm{s^2}$, $[\gamma_j]=\mathrm{s^{-1}}$ and
$[\omega_j]=\mathrm{s^{-1}}$. Each oscillator has supplied
$\mathbf P_j(t_0),\mathbf V_j(t_0)$. This is a selected Lorentz material model,
not a universal constitutive law. Its optical origin is described in
\href{https://www.ocw.mit.edu/courses/6-974-fundamentals-of-photonics-quantum-electronics-spring-2006/8e54d1625e9e71eb5b4e4998e6967e4e_chapter2.pdf}{MIT Fundamentals of Photonics, Section 2.1.4 [EM4]}.

\begin{proposition}[Explicit storage and dissipation]
For the held coefficients in equation~\eqref{eq:r13-em-lorentz}, define
\begin{align}
 u_j&=\frac{|\mathbf V_j|^2+\omega_j^2|\mathbf P_j|^2}{2f_j},\nonumber\\
 u_*&=\tfrac12\epsilon_b|\mathbf E|^2+
       \tfrac12\mu|\mathbf H|^2+\sum_j u_j,\nonumber\\
 q_{\rm diss}&=\mathbf E\cdot\sigma\mathbf E+
                    \sum_j\frac{\gamma_j}{f_j}|\mathbf V_j|^2\ge0.
 \label{eq:r13-em-storage}
\end{align}
Then $\partial_tu_*+\nabla\cdot\mathbf S
=-\mathbf E\cdot\mathbf J_{\rm src}-q_{\rm diss}$.
\end{proposition}
\begin{proof}
Multiplying the oscillator equation by $\mathbf V_j/f_j$ gives
$\dot u_j=\mathbf E\cdot\mathbf V_j-
\gamma_j|\mathbf V_j|^2/f_j$.
In equation~\eqref{eq:r13-em-poynting-general}, the polarization part of
$\mathbf E\cdot\dot{\mathbf D}$ is $\sum_j\mathbf E\cdot\mathbf V_j$.
Adding the oscillator balances cancels these exchanges exactly. Positivity
follows from the coefficient and conductivity hypotheses.
\end{proof}

This supplies a consistent source for R12's thermal ledger. Assign
$P_{\rm heat,opt}=\int_\Omega q_{\rm diss}\,dV$ only when these dissipative
channels thermalize in the modeled thermal store on the declared interval.
If excitation, emission or another store participates, partition the sink and
retain that store's evolution. Counting the same absorbed power independently
as full heat and full chemical excitation would break the balance.

With the phasor convention $\mathbf E(t)=\operatorname{Re}
(\widehat{\mathbf E}e^{-i\omega t})$, $\omega>0$, the scalar isotropic,
steady harmonic closure is
\begin{align}
 \epsilon_{\rm eff}(\omega)&=\epsilon_b+
 \sum_j\frac{f_j}{\omega_j^2-\omega^2-i\gamma_j\omega}
 +\frac{i\sigma}{\omega},\nonumber\\
 \langle q_{\rm diss}\rangle&=
 \frac{\omega}{2}\operatorname{Im}(\epsilon_{\rm eff})
                  |\widehat{\mathbf E}|^2.
 \label{eq:r13-em-epsilon}
\end{align}
All denominators must be nonzero, and the selected steady response must exist;
an undamped resonance does not satisfy these assumptions. Conductivity is
included once, inside $\epsilon_{\rm eff}$ in this frequency convention.
The time-domain state in equation~\eqref{eq:r13-em-lorentz} remains the
definition when a steady response is inappropriate.

Temperature, hysteresis and reporter state may parameterize a family of these
material profiles. Holding them fixed during an optical subproblem is an
explicit time-scale approximation. If coefficients change during propagation,
their derivatives change stored energy: for example,
$\partial_t(\epsilon_b|\mathbf E|^2/2)$ includes
$\dot\epsilon_b|\mathbf E|^2/2$. The corresponding modulation work and
material exchange must be included. Replacing a coefficient at each tick
without that accounting is not the proved coupled energy law.

\subsection{How Maxwell closes the existing propagation operator}
\label{sec:r13-em-propagation}

In a homogeneous, source-free, lossless region with constant scalar
$\epsilon,\mu>0$, the divergence constraint gives $\nabla\cdot\mathbf E=0$.
Taking a curl and using $\nabla\times\nabla\times\mathbf E
=\nabla(\nabla\cdot\mathbf E)-\nabla^2\mathbf E$ yields
\begin{align}
 \nabla^2\mathbf E-\mu\epsilon\partial_t^2\mathbf E&=0,\nonumber\\
 (\nabla^2+k^2)\widehat{\mathbf E}&=0,
 \quad k=\omega\sqrt{\mu\epsilon}.
 \label{eq:r13-em-wave}
\end{align}
For a linear dispersive homogeneous medium, use the corresponding frequency
equation with its declared complex constitutive coefficients and radiation
branch. A scalar component can stand for the optical field only within an
appropriate polarization/geometry reduction. Spatially varying tensors,
strong polarization conversion or near-field boundary interactions do not
automatically reduce to the same scalar equation.
\href{https://ocw.mit.edu/courses/2-71-optics-spring-2014/e6cfb73b15ef5b0610fe662794eb3205_MIT2_71S14_lec9_notes.pdf}{MIT Optics, Lecture 9 [EM5]}.

For the scalar, forward, homogeneous lossless case write $U=e^{ikz}\psi$.
Substitution gives
\begin{equation}
 \partial_z^2\psi+2ik\partial_z\psi+\nabla_\perp^2\psi=0.
 \label{eq:r13-em-paraxial-exact}
\end{equation}
The paraxial approximation drops $\partial_z^2\psi$ when it is negligible
relative to the retained terms on the admitted field class. For $z>0$, one
resulting Fresnel transfer, with $\lambda_m=2\pi/k$, is
\begin{equation}
 (H_zU_0)(x,y)=\frac{e^{ikz}}{i\lambda_m z}
 \iint U_0(x',y')
 \exp\!\left(\frac{ik[(x-x')^2+(y-y')^2]}{2z}\right)dx'\,dy'.
 \label{eq:r13-em-fresnel}
\end{equation}
This states the normalization missing from a mere angular-shape transform.
For R12's finite aperture, the integral is over its bounded support; its exact
real-pair lift needs the section integrability conditions already specified
there. A full-plane operator has its separately declared function-space
domain. Fraunhofer additionally neglects the aperture-plane quadratic phase
on the stated field/support scale. These approximations need application
error allowances, not an assertion that exact encoding makes them exact
Maxwell solutions.
\href{https://ocw.mit.edu/courses/2-71-optics-spring-2014/4b53a5747f9b58b9b73b2bbcc39945f0_MIT2_71S14_lec16_notes.pdf}{MIT Optics, Lecture 16 [EM6]}.

The same boundary-value problem fixes $t_\lambda$ and reflected/outgoing
channels from material parameters and geometry. A thin local multiplier is
an adopted reduction; a general scattering device may require a matrix or
integral transfer. The aperture $A$ remains a geometric predicate. It does
not specify the polarization currents or electromagnetic boundary response
of the physical opaque material by itself.

\subsection{Passivity at the same physical interface}
\label{sec:r13-em-passivity}

For a fixed volume, let $\mathcal U=\int_\Omega u_*dV$. Separate its boundary
flux into admitted incoming and outgoing power. With no impressed source,
integration of the derived energy identity gives
\begin{equation}
 E_{\rm out}[t_0,t_1]\le E_{\rm in}[t_0,t_1]
                   +\mathcal U(t_0)-\mathcal U(t_1).
 \label{eq:r13-em-transient-passive}
\end{equation}
Stored energy can temporarily power an output. In steady harmonic operation,
the average storage change vanishes; for independent positive-flux modal
coordinates this gives precisely R12's condition
\begin{equation}
 S^\dagger W_{\rm out}S\preceq W_{\rm in}.
 \label{eq:r13-em-passive-matrix}
\end{equation}
Both metrics refer to the same physical normalization and selected ports.
Unobserved reflected, radiated or coupled-away power is not automatically
absorption. Only with all outgoing channels accounted for, no supplied gain
and no net storage accumulation does $\mathcal R+\mathcal T+\mathcal A=1$
identify the remaining fraction as material absorption. This derivation
supplies the physical premise for the R12 scaled passive pinion construction;
it does not change its algebra or demonstrate a manufactured device.

One optional mechanical coupling is already useful without a general
electromagnetic stress model. For a stationary planar target in vacuum,
normal incidence, forward transmission and specular backward reflection,
the steady force along the incident beam is
\begin{equation}
 F_\parallel=\frac{P_{\rm in}}{c}(1+\mathcal R-\mathcal T)
            =\frac{P_{\rm in}}{c}(\mathcal A+2\mathcal R).
 \label{eq:r13-em-pressure}
\end{equation}
It follows by subtracting outgoing from incoming momentum flux. It can enter
the declared mechanical load. Oblique/diffuse scattering, moving media,
directional re-emission or additional forces require their corresponding
momentum balance; the scalar formula does not cover them.
\href{https://www.feynmanlectures.caltech.edu/I_34.html}{Caltech Feynman Lectures, I.34.9 [EM7]}.

\subsection{Finite exact seed representation}
\label{sec:r13-em-finite}

Complex coefficients, amplitudes and matrices use R12's ordered real-pair
templates. For a declared finite modal problem $Az=b$, with $A=X+iY$,
$z=p+iq$ and $b=r+is$, the literal real system is
\begin{equation}
 \begin{pmatrix}X&-Y\\Y&X\end{pmatrix}
 \begin{pmatrix}p\\q\end{pmatrix}=
 \begin{pmatrix}r\\s\end{pmatrix}.
 \label{eq:r13-em-real-system}
\end{equation}
An inverse requires a proved nonzero determinant; an unresolved singularity
or branch remains an explicit non-value outcome. The coefficients bind the
chosen basis, geometry, boundary projection, material law and frequency.
They may be retained exact expressions or supplied calibrated literals with
uncertainty. For instance, select a bounded fixed cavity with a perfect electric
conductor boundary and real admissible basis functions $\boldsymbol\psi_j$
whose tangential traces vanish. The expansion
$\widehat{\mathbf E}_N=\sum_{j=1}^{N}z_j\boldsymbol\psi_j$ and weak projection
of the frequency curl--curl equation give
\begin{align}
 A_{ab}&=\int_\Omega\!\left[
 (\nabla\times\boldsymbol\psi_a)\cdot\mu^{-1}
 (\nabla\times\boldsymbol\psi_b)
 -\omega^2\boldsymbol\psi_a\cdot\epsilon_{\rm eff}
             \boldsymbol\psi_b\right]dV,\nonumber\\
 b_a&=i\omega\int_\Omega\boldsymbol\psi_a\cdot
                    \widehat{\mathbf J}_{\rm src}\,dV.
 \label{eq:r13-em-modal}
\end{align}
The boundary term vanishes under this particular trace condition. The finite
space, source projection, divergence constraints and coefficient integral
domains must be specified. Other boundaries require their own terms. These
volume-integral definitions fit existing exact operators only when their
actual integrands and bounded coordinate domains meet that operator contract;
the weak notation alone is not an executable integral encoding.

The finite polarization coordinates use the first-order equations in
\eqref{eq:r13-em-lorentz}. Together with a declared finite field reduction,
they form an ordinary finite state evolution only after building algebraic
constraints into the basis or explicitly eliminating them to obtain a
nonsingular finite evolution. The fixed representation, selected elimination
branch and regular domain remain explicit. XOP can denote the resulting ODE
under its existence/domain contract or an explicitly selected finite update.
Remaining algebraic constraints need their own defined constrained formulation.
It does not acquire an arbitrary continuum PDE solver. Basis truncation,
omitted modes and an approximate constitutive fit remain model errors even
when every retained operation and packed bit is exact. The physical profile
therefore exports the field transfer, dissipative/other energy channels and
optional force from one common model, rather than fitting unrelated outputs
that happen to share a seed identifier.
